\documentclass[10pt,a4paper]{article}
\usepackage[utf8]{inputenc}
\usepackage[vietnamese]{babel}
\usepackage{amsmath,amssymb,amsfonts}
\usepackage{geometry}
\geometry{a4paper, left=10mm, right=10mm, top=10mm, bottom=10mm}
\usepackage{tikz}
\usetikzlibrary{shapes,arrows.meta,positioning,calc}
\usepackage[most]{tcolorbox}
\usepackage{array,booktabs,multicol}
\usepackage{fontawesome}
\usepackage{helvet}
\renewcommand{\familydefault}{\sfdefault}
\usepackage{parskip}

\definecolor{primary}{RGB}{112, 48, 160}
\definecolor{teal}{RGB}{0, 128, 128}
\definecolor{orange}{RGB}{230, 108, 44}
\definecolor{green}{RGB}{46, 139, 87}
\definecolor{rose}{RGB}{192, 0, 0}

\pagestyle{empty}

\begin{document}

% --- HEADER BANNER ---
\begin{tcolorbox}[
  enhanced,
  colback=violet!8!white,
  colframe=violet!80!black,
  arc=2.5mm, boxrule=1pt,
  left=3mm, right=3mm, top=2mm, bottom=2mm
]
  \begin{tabular*}{\textwidth}{@{\extracolsep{\fill}} lr}
    {\small\bfseries\color{violet!85!black}\faBook\ SỔ TAY KIẾN THỨC TOÁN 2} & 
    {\small\bfseries\color{teal!85!black}\faGraduationCap\ CHƯƠNG TRÌNH GDPT 2018}
  \end{tabular*}
  \vspace{-1mm}
  \begin{center}
    \rule{0.98\textwidth}{0.6pt}
  \end{center}
  \vspace{-2mm}
  \centering
  {\LARGE\bfseries\color{violet!90!black}\faBookmark\ BÀI 56: GIỚI THIỆU TIỀN VIỆT NAM}\\[1mm]
  {\itshape\footnotesize\color{black!75} Nhận biết các tờ tiền $100$ đồng, $200$ đồng, $500$ đồng, $1\,000$ đồng}
\end{tcolorbox}

\vspace{-2mm}

\noindent
\begin{minipage}[t]{0.49\textwidth}\vspace{0pt}
\begin{tcolorbox}[
  enhanced,
  colback=violet!4!white,
  colframe=violet!75!black,
  arc=2mm, boxrule=0.8pt,
  left=2mm, right=2mm, top=1.8mm, bottom=1.8mm,
  title={\footnotesize\bfseries\color{white}\faKey\ 1. Nhận biết các mệnh giá tiền Việt Nam}
]
{\footnotesize
\begin{itemize}\setlength{\itemsep}{1pt}\setlength{\leftskip}{-2.5mm}
  \item Đơn vị tiền tệ của nước ta là \textbf{đồng}.
  \item Các tờ tiền giấy trong phạm vi học tập lớp 2:
  \begin{itemize}\setlength{\itemsep}{0.5pt}\setlength{\leftskip}{-2mm}
    \item Tờ \textbf{$100$ đồng} (một trăm đồng).
    \item Tờ \textbf{$200$ đồng} (hai trăm đồng).
    \item Tờ \textbf{$500$ đồng} (năm trăm đồng).
    \item Tờ \textbf{$1\,000$ đồng} (một nghìn đồng).
  \end{itemize}
  \item Trên mỗi tờ tiền đều ghi:
  \begin{itemize}\setlength{\itemsep}{0.5pt}\setlength{\leftskip}{-2mm}
    \item Dòng chữ: \textit{"CỘNG HÒA XÃ HỘI CHỦ NGHĨA VIỆT NAM"}.
    \item Mệnh giá bằng \textbf{số} và bằng \textbf{chữ}.
    \item Chân dung Chủ tịch Hồ Chí Minh và Quốc huy.
  \end{itemize}
\end{itemize}
}
\end{tcolorbox}
\end{minipage}
\hfill
\begin{minipage}[t]{0.49\textwidth}\vspace{0pt}
\begin{tcolorbox}[
  enhanced,
  colback=teal!4!white,
  colframe=teal!75!black,
  arc=2mm, boxrule=0.8pt,
  left=2mm, right=2mm, top=1.8mm, bottom=1.8mm,
  title={\footnotesize\bfseries\color{white}\faSliders\ 2. Kĩ năng đếm tiền và đổi tiền}
]
{\footnotesize
\begin{itemize}\setlength{\itemsep}{1pt}\setlength{\leftskip}{-2.5mm}
  \item \textbf{Đếm tổng số tiền:} Cộng mệnh giá của từng tờ tiền lại với nhau.
  \item \textbf{Mối quan hệ đổi tiền:}
  \begin{itemize}\setlength{\itemsep}{0.5pt}\setlength{\leftskip}{-2mm}
    \item $1$ tờ $1\,000$ đồng $= 2$ tờ $500$ đồng ($500 + 500 = 1\,000$).
    \item $1$ tờ $1\,000$ đồng $= 5$ tờ $200$ đồng.
    \item $1$ tờ $1\,000$ đồng $= 10$ tờ $100$ đồng.
    \item $1$ tờ $500$ đồng $= 2$ tờ $200$ đồng $+ 1$ tờ $100$ đồng.
  \end{itemize}
  \item Kĩ năng tính tiền thừa khi mua sắm:
  $$\text{Tiền thừa} = \text{Tiền đưa} - \text{Giá tiền món hàng}$$
\end{itemize}
}
\end{tcolorbox}
\end{minipage}

\vspace{-1.5mm}
\begin{tcolorbox}[
  enhanced,
  colback=orange!4!white,
  colframe=orange!80!black,
  arc=2mm, boxrule=0.8pt,
  left=2mm, right=2mm, top=1.5mm, bottom=1.5mm,
  title={\footnotesize\bfseries\color{white}\faSitemap\ 3. Sơ đồ các tờ tiền và mối quan hệ quy đổi tương đương}
]
\centering
\begin{tikzpicture}[x=1.35cm, y=0.75cm]
  \tikzset{arr/.style={-{Stealth[scale=0.8]}}, every node/.append style={align=center}}
  \node[rectangle, rounded corners=2pt, draw=red!80!black, fill=red!15, font=\footnotesize\bfseries] (k) at (0,0) {$\mathbf{1\,000}$ ĐỒNG};
  \node[rectangle, rounded corners=2pt, draw=violet!80!black, fill=violet!15, font=\scriptsize\bfseries] (c5) at (3.2,0.8) {$2$ tờ $\mathbf{500}$ đồng\\($500 + 500 = 1\,000$)};
  \node[rectangle, rounded corners=2pt, draw=teal!80!black, fill=teal!15, font=\scriptsize\bfseries] (c2) at (3.2,0) {$5$ tờ $\mathbf{200}$ đồng\\($200 \times 5 = 1\,000$)};
  \node[rectangle, rounded corners=2pt, draw=orange!80!black, fill=orange!15, font=\scriptsize\bfseries] (c1) at (3.2,-0.8) {$10$ tờ $\mathbf{100}$ đồng\\($100 \times 10 = 1\,000$)};

  \draw[arr, thick] (k) |- (c5);
  \draw[arr, thick] (k) -- (c2);
  \draw[arr, thick] (k) |- (c1);
\end{tikzpicture}
\end{tcolorbox}

\vspace{-1.5mm}
\noindent
\begin{minipage}[t]{0.49\textwidth}\vspace{0pt}
\begin{tcolorbox}[
  enhanced,
  colback=green!4!white,
  colframe=green!60!black,
  arc=2mm, boxrule=0.8pt,
  left=2mm, right=2mm, top=1.8mm, bottom=1.8mm,
  title={\footnotesize\bfseries\color{white}\faCheckCircle\ 4. Ví dụ thực hành mua bán}
]
{\footnotesize
\textbf{1. Tính tổng số tiền:}\\
Nam có $1$ tờ $500$ đồng và $2$ tờ $200$ đồng.\\
$\implies$ Tổng tiền $= 500 + 200 + 200 = \mathbf{900}$ đồng.\\
\textbf{2. Bài toán trả lại tiền thừa:}\\
Hà có $1\,000$ đồng. Hà mua một chiếc kẹo giá $700$ đồng. Hỏi cô bán hàng trả lại Hà bao nhiêu tiền?\\
\textit{Bài giải:}\\
Số tiền cô bán hàng trả lại Hà là:\\
$1\,000 - 700 = \mathbf{300}$ (đồng).
}
\end{tcolorbox}
\end{minipage}
\hfill
\begin{minipage}[t]{0.49\textwidth}\vspace{0pt}
\begin{tcolorbox}[
  enhanced,
  colback=rose!4!white,
  colframe=rose!70!black,
  arc=2mm, boxrule=0.8pt,
  left=2mm, right=2mm, top=1.8mm, bottom=1.8mm,
  title={\footnotesize\bfseries\color{white}\faExclamationTriangle\ 5. Lời dặn dò thực tế}
]
{\footnotesize
\begin{itemize}\setlength{\itemsep}{1.5pt}\setlength{\leftskip}{-2.5mm}
  \item \textbf{Luôn kèm chữ "đồng"} vào sau kết quả số tiền (ví dụ: $800$ đồng, không viết trơ trọi $800$).
  \item Giữ gìn tiền phẳng phiu, sạch sẽ, không vẽ bậy, không làm rách hay nhàu nát tiền.
  \item Biết tiết kiệm tiền bạc, chi tiêu hợp lí, không vòi vĩnh bố mẹ mua sắm lãng phí.
\end{itemize}
}
\end{tcolorbox}
\end{minipage}


\end{document}
